
\section{Rocket League}

O trabalho utiliza \textit{Rocket League} como primeiro caso de análise. O jogo apresenta veículos e ambientes com materiais e texturas que contribuem para a aparência dos objetos. A documentação consultada para o trabalho inclui material da Substance 3D relacionado ao processo de modelagem e texturização de Rocket League.

No contexto deste trabalho, a técnica é relevante porque as texturas precisam ser associadas às superfícies dos veículos e elementos do cenário. O UV mapping é o mecanismo que estabelece essa correspondência entre o espaço da textura e a superfície 3D.

A análise visual deve ser complementada pela fonte técnica utilizada, pois uma captura de gameplay mostra o resultado final, mas não permite observar diretamente as coordenadas UV internas ou o shader utilizado.

\section{Marvel's Spider-Man Remastered}

O segundo caso analisado é \textit{Marvel's Spider-Man Remastered}. O jogo utiliza modelos tridimensionais, materiais, texturas e shaders para produzir a aparência dos personagens e ambientes. O material da Insomniac Games citado na pesquisa apresenta diferenças gráficas da versão remasterizada, enquanto materiais de modificação também discutem aspectos relacionados a texturas e UVs.

Nesse caso, o UV mapping participa da associação de texturas às superfícies dos modelos, incluindo detalhes visuais das roupas e demais elementos. Assim como no primeiro caso, a captura de gameplay serve como evidência visual, enquanto as fontes técnicas sustentam a interpretação.

\section{Comparação}

\begin{table}[H]
\centering
\caption{Comparação dos casos analisados.}
\begin{tabular}{p{3.2cm}p{5cm}p{5cm}}
\toprule
\textbf{Jogo} & \textbf{Elemento observado} & \textbf{Relação com a técnica} \\
\midrule
Rocket League & Veículos e elementos do estádio & Associação de texturas e materiais às superfícies 3D. \\
Marvel's Spider-Man Remastered & Trajes, personagens e ambientes & Uso de texturas, materiais e shaders associados aos modelos. \\
\bottomrule
\end{tabular}
\end{table}
